\hypertarget{python-3.0-the-unicode-paradigm-shift-and-text-vs.-bytes-separation}{%
\chapter{Python 3.0: The Unicode Paradigm Shift and Text vs.~Bytes
Separation}\label{python-3.0-the-unicode-paradigm-shift-and-text-vs.-bytes-separation}}

\hypertarget{the-unicode-paradigm-shift-and-text-vs.-bytes-separation-pep-358-pep-3112}{%
\subsection{7.1 The Unicode Paradigm Shift and Text vs.~Bytes Separation
(PEP 358 \& PEP
3112)}\label{the-unicode-paradigm-shift-and-text-vs.-bytes-separation-pep-358-pep-3112}}

Python 3.0 introduced a structural boundary between textual characters
and binary data. In Python 2.x, \passthrough{\lstinline!str!} served
double-duty as both raw bytes and text, causing silent,
non-deterministic bugs when ASCII decoding failed implicitly during
concatenation. Python 3.0 solved this by establishing a strict boundary
between \passthrough{\lstinline!str!} (text) and
\passthrough{\lstinline!bytes!}/\passthrough{\lstinline!bytearray!}
(binary).

\hypertarget{the-python-2.x-concatenation-trap-and-coercion-ban}{%
\subsubsection{1. The Python 2.x Concatenation Trap and Coercion
Ban}\label{the-python-2.x-concatenation-trap-and-coercion-ban}}

In Python 2.x, mixed operations between \passthrough{\lstinline!str!}
(raw 8-bit characters) and \passthrough{\lstinline!unicode!} (arbitrary
code points) were implicitly resolved. When executing
\passthrough{\lstinline!"rsum" + u" (French)"!}, CPython attempted to
promote the \passthrough{\lstinline!str!} by decoding it via the default
codec (usually ASCII):

\begin{lstlisting}[language=C]
coerced = PyUnicode_FromEncodedObject(str_obj, "ascii", "strict");
\end{lstlisting}

If the raw string contained non-ASCII bytes (e.g.,
\passthrough{\lstinline!\\xe9!} for Latin-1, or
\passthrough{\lstinline!\\xc3\\xa9!} for UTF-8), this implicit decoding
step raised a \passthrough{\lstinline!UnicodeDecodeError!} at runtime.
Python 3.0 eliminated this by raising an unconditional
\passthrough{\lstinline!TypeError!} on any implicit mixed-type
operation:

\begin{lstlisting}[language=Python]
# Python 3.0 runtime coercion ban
>>> b"data" + "string"
Traceback (most recent call last):
  File "<stdin>", line 1, in <module>
TypeError: can't concat bytes to str
\end{lstlisting}

At the C-API level, binary and text domains are kept completely
disjoint. \passthrough{\lstinline!PyUnicode\_Compare!} returns a failure
state if one operand is a bytes object, and
\passthrough{\lstinline!PyBytes\_Concat!} does not accept unicode
inputs.

\hypertarget{c-level-representations-of-binary-objects-pybytesobject-pybytearrayobject}{%
\subsubsection{\texorpdfstring{2. C-Level Representations of Binary
Objects: \texttt{PyBytesObject} \&
\texttt{PyByteArrayObject}}{2. C-Level Representations of Binary Objects: PyBytesObject \& PyByteArrayObject}}\label{c-level-representations-of-binary-objects-pybytesobject-pybytearrayobject}}

CPython represents the immutable \passthrough{\lstinline!bytes!} type at
the C-level as \passthrough{\lstinline!PyBytesObject!}, defined in
\passthrough{\lstinline!Include/bytesobject.h!}:

\begin{lstlisting}[language=C]
typedef struct {
    PyObject_VAR_HEAD
    Py_hash_t ob_shash;         /* Cached hash value of bytes; -1 if uncomputed */
    char ob_sval[1];            /* Contiguous byte vector; dynamically sized to [ob_size + 1] */
} PyBytesObject;
\end{lstlisting}

Here, \passthrough{\lstinline!ob\_sval!} is a character array that is
allocated inline with the rest of the object structure, terminating with
a trailing \passthrough{\lstinline!\\0!} to remain compatible with
standard C library functions (e.g., \passthrough{\lstinline!strlen!} and
\passthrough{\lstinline!strcpy!}).

The mutable sibling, \passthrough{\lstinline!bytearray!}, is represented
by \passthrough{\lstinline!PyByteArrayObject!} defined in
\passthrough{\lstinline!Include/bytearrayobject.h!}:

\begin{lstlisting}[language=C]
typedef struct {
    PyObject_VAR_HEAD
    Py_ssize_t ob_alloc;        /* Number of bytes allocated in the buffer */
    char *ob_bytes;             /* Pointer to the start of the allocated memory block */
    char *ob_start;             /* Pointer to the start of the logical byte sequence */
    Py_ssize_t ob_exports;      /* Reference count of active buffer exports (locking) */
} PyByteArrayObject;
\end{lstlisting}

\passthrough{\lstinline!ob\_bytes!} and
\passthrough{\lstinline!ob\_start!} are split to allow efficient
head-deletion/prepends. When a user calls
\passthrough{\lstinline!del array[0]!},
\passthrough{\lstinline!ob\_start!} is simply incremented and
\passthrough{\lstinline!ob\_size!} decremented, avoiding an \(O(n)\)
memory shift. Memory buffer growth is managed via an overallocation
growth factor:
\[\text{ob\_alloc} = \text{new\_size} + (\text{new\_size} \gg 3) + (\text{new\_size} < 9 \,?\, 3 : 6)\]
This guarantees that append operations have an amortized \(O(1)\) time
complexity. \passthrough{\lstinline!ob\_exports!} tracks references from
active \passthrough{\lstinline!memoryview!} objects; if
\passthrough{\lstinline!ob\_exports > 0!}, any operation that attempts
to resize or reallocate the underlying
\passthrough{\lstinline!ob\_bytes!} buffer will raise a
\passthrough{\lstinline!BufferError!}.

\hypertarget{low-level-c-representation-pre-pep-393-pyunicodeobject}{%
\subsubsection{\texorpdfstring{3. Low-Level C Representation: Pre-PEP
393
\texttt{PyUnicodeObject}}{3. Low-Level C Representation: Pre-PEP 393 PyUnicodeObject}}\label{low-level-c-representation-pre-pep-393-pyunicodeobject}}

Prior to the string optimization introduced in Python 3.3 (PEP 393),
CPython represented Unicode strings using a uniform array of
\passthrough{\lstinline!Py\_UNICODE!} units.
\passthrough{\lstinline!PyUnicodeObject!} was defined in
\passthrough{\lstinline!Include/unicodeobject.h!} as:

\begin{lstlisting}[language=C]
typedef struct {
    PyObject_HEAD
    Py_ssize_t length;          /* Number of code points */
    Py_UNICODE *str;            /* Pointer to the character array */
    Py_hash_t hash;             /* Cached hash value; -1 if uncomputed */
    PyObject *defenc;           /* Cached UTF-8 encoded bytes representation */
} PyUnicodeObject;
\end{lstlisting}

The representation type \passthrough{\lstinline!Py\_UNICODE!} was
defined at compile-time as: * \textbf{UCS-2 (Narrow Build)}: Compiled
with \passthrough{\lstinline!wchar\_t!} as a 16-bit type (size 2 bytes).
Characters beyond \passthrough{\lstinline!U+FFFF!} (e.g., emojis) had to
be represented as surrogate pairs, causing index and len calculations to
mismatch. * \textbf{UCS-4 (Wide Build)}: Compiled with
\passthrough{\lstinline!wchar\_t!} as a 32-bit type (size 4 bytes).
Every Unicode code point mapped to exactly one array index, but
ASCII-only strings consumed 4 times the memory they required.

\hypertarget{pep-383-the-surrogateescape-error-handler-mechanics}{%
\subsubsection{\texorpdfstring{4. PEP 383: The \texttt{surrogateescape}
Error Handler
Mechanics}{4. PEP 383: The surrogateescape Error Handler Mechanics}}\label{pep-383-the-surrogateescape-error-handler-mechanics}}

To allow POSIX filesystems (which use arbitrary null-terminated byte
sequences for paths) to round-trip pathnames containing invalid UTF-8
bytes without throwing exceptions, PEP 383 introduced the
\passthrough{\lstinline!surrogateescape!} error handler. During
decoding, any invalid byte (which does not form a valid UTF-8 sequence)
is mapped to a high-surrogate code point in the range
\passthrough{\lstinline!U+DC80!} to \passthrough{\lstinline!U+DCFF!}
via: \[\text{code\_point} = 0\text{xDC00} + \text{byte\_value}\] During
encoding, these specific surrogate code points are mapped back to their
original raw bytes:
\[\text{byte\_value} = \text{code\_point} - 0\text{xDC00}\] This allows
arbitrary binary data to be round-tripped through Unicode
\passthrough{\lstinline!str!} representations:

\begin{lstlisting}[language=Python]
# Raw undecodable path round-trip simulation
raw_bytes = b"bad_\xff_path.txt"
decoded_str = raw_bytes.decode("utf-8", "surrogateescape")
# decoded_str contains code point U+DCFF where \xff was.
reencoded_bytes = decoded_str.encode("utf-8", "surrogateescape")
assert reencoded_bytes == raw_bytes
\end{lstlisting}

The CPython implementation in
\passthrough{\lstinline!Objects/unicodeobject.c!} handles this matching
logic:

\begin{lstlisting}[language=C]
/* Pseudocode of surrogateescape decoding branch */
if (status == INVALID_BYTE) {
    *unicode_ptr++ = 0xDC00 + (unsigned char)input_byte;
}
\end{lstlisting}

\begin{center}\rule{0.5\linewidth}{0.5pt}\end{center}

\hypertarget{the-print-function-redesign}{%
\subsection{\texorpdfstring{7.2 The \texttt{print()} Function
Redesign}{7.2 The print() Function Redesign}}\label{the-print-function-redesign}}

\hypertarget{bytecode-comparison}{%
\subsubsection{1. Bytecode Comparison}\label{bytecode-comparison}}

In Python 2.x, \passthrough{\lstinline!print!} was a core language
statement with specialized bytecodes. In Python 3.0, it was unified into
a standard built-in function.

\hypertarget{python-2.7-compiler-translation-for-print-hello}{%
\paragraph{\texorpdfstring{Python 2.7 compiler translation for
\texttt{print\ "hello"}:}{Python 2.7 compiler translation for print "hello":}}\label{python-2.7-compiler-translation-for-print-hello}}

\begin{lstlisting}
1           0 LOAD_CONST               0 ('hello')
            3 PRINT_ITEM
            4 PRINT_NEWLINE
\end{lstlisting}

The VM execution loop (\passthrough{\lstinline!ceval.c!}) maps
\passthrough{\lstinline!PRINT\_ITEM!} directly to standard output stream
writing operations, meaning the behavior was fixed at compilation.

\hypertarget{python-3.0-compiler-translation-for-printhello}{%
\paragraph{\texorpdfstring{Python 3.0 compiler translation for
\texttt{print("hello")}:}{Python 3.0 compiler translation for print("hello"):}}\label{python-3.0-compiler-translation-for-printhello}}

\begin{lstlisting}
1           0 LOAD_GLOBAL              0 (print)
            3 LOAD_CONST               1 ('hello')
            6 CALL_FUNCTION            1
\end{lstlisting}

The interpreter dynamically resolves \passthrough{\lstinline!print!} at
runtime via standard namespace lookups. This enables runtime overriding:

\begin{lstlisting}[language=Python]
import builtins
def custom_print(*args, **kwargs):
    builtins.print("[LOG]", *args, **kwargs)
builtins.print = custom_print
\end{lstlisting}

\hypertarget{c-level-stream-writing-textiowrapper}{%
\subsubsection{2. C-Level stream writing \&
TextIOWrapper}\label{c-level-stream-writing-textiowrapper}}

The C-level entrypoint for the \passthrough{\lstinline!print()!}
function is \passthrough{\lstinline!builtin\_print!} inside
\passthrough{\lstinline!Python/bltinmodule.c!}:

\begin{lstlisting}[language=C]
static PyObject *
builtin_print(PyObject *self, PyObject *args, PyObject *kwds)
{
    static char *kwlist[] = {"sep", "end", "file", "flush", 0};
    PyObject *sep = Py_None, *end = Py_None, *file = Py_None;
    int flush = 0;
    // Parsing keywords using PyArg_ParseTupleAndKeywords...
    ...
}
\end{lstlisting}

\begin{itemize}
\tightlist
\item
  \textbf{Stream Resolution}: If \passthrough{\lstinline!file!} is
  \passthrough{\lstinline!Py\_None!} or omitted, the function queries
  \passthrough{\lstinline!sys.stdout!} via
  \passthrough{\lstinline!PySys\_GetObject("stdout")!}.
\item
  \textbf{String Conversions}: For each positional argument,
  \passthrough{\lstinline!builtin\_print!} calls
  \passthrough{\lstinline!PyObject\_Str(arg)!} to compute its string
  representation.
\item
  \textbf{Buffer Flushing}: After calling
  \passthrough{\lstinline!file.write()!}, if
  \passthrough{\lstinline!flush=True!} is set, CPython attempts to
  invoke the \passthrough{\lstinline!flush()!} method of the
  \passthrough{\lstinline!file!} object. If
  \passthrough{\lstinline!sys.stdout!} wraps a standard output
  descriptor (stdout), this triggers
  \passthrough{\lstinline!fflush(stdout)!} in the underlying C library,
  executing a synchronous write syscall.
\end{itemize}

\begin{center}\rule{0.5\linewidth}{0.5pt}\end{center}

\hypertarget{division-unification-pep-238}{%
\subsection{7.3 Division Unification (PEP
238)}\label{division-unification-pep-238}}

\hypertarget{the-division-operators-and-bytecodes}{%
\subsubsection{1. The Division Operators and
Bytecodes}\label{the-division-operators-and-bytecodes}}

Python 2.x's division operator \passthrough{\lstinline!/!} mapped to the
\passthrough{\lstinline!BINARY\_DIVIDE!} bytecode, which executed floor
division if both operands were integers, and true division if either was
a float. This led to silent precision loss bugs. Python 3.0 unified this
by assigning \passthrough{\lstinline!/!} to true division and
introducing \passthrough{\lstinline!//!} for floor division: *
\textbf{True Division (\passthrough{\lstinline!/!})} maps to the
\passthrough{\lstinline!BINARY\_TRUE\_DIVIDE!} bytecode, which always
delegates to the \passthrough{\lstinline!nb\_true\_divide!} slot. *
\textbf{Floor Division (\passthrough{\lstinline!//!})} maps to the
\passthrough{\lstinline!BINARY\_FLOOR\_DIVIDE!} bytecode, delegating to
the \passthrough{\lstinline!nb\_floor\_divide!} slot.

\hypertarget{c-level-integer-division-mechanics}{%
\subsubsection{2. C-Level Integer Division
Mechanics}\label{c-level-integer-division-mechanics}}

CPython defines number methods in the
\passthrough{\lstinline!PyNumberMethods!} struct inside
\passthrough{\lstinline!Include/object.h!}:

\begin{lstlisting}[language=C]
typedef struct {
    binaryfunc nb_add;
    binaryfunc nb_subtract;
    ...
    binaryfunc nb_floor_divide;
    binaryfunc nb_true_divide;
} PyNumberMethods;
\end{lstlisting}

For integer types, these slots point to C functions in
\passthrough{\lstinline!Objects/longobject.c!}: *
\passthrough{\lstinline!nb\_floor\_divide!} points to
\passthrough{\lstinline!long\_div!}. It computes the integer quotient
and rounds towards negative infinity. *
\passthrough{\lstinline!nb\_true\_divide!} points to
\passthrough{\lstinline!long\_true\_divide!}.

The execution path of
\passthrough{\lstinline!long\_true\_divide(x, y)!}: 1. It extracts the
size of integers \(x\) and \(y\). If they fit within C double precision
(53 bits of mantissa), it converts them to double:
\[x_{\text{double}} = (\text{double})x; \quad y_{\text{double}} = (\text{double})y;\]
2. If the integer values exceed double limits, CPython runs an
arbitrary-precision float conversion algorithm
(\passthrough{\lstinline!\_PyLong\_Format!} or manual bit shift scaling)
to extract the most significant 53 bits. 3. It performs the float
division and instantiates a new \passthrough{\lstinline!PyFloatObject!}
to hold the output:
\[\text{result} = x_{\text{double}} / y_{\text{double}}\] 4. If
\(y = 0\), it raises a \passthrough{\lstinline!ZeroDivisionError!} via
\passthrough{\lstinline!PyErr\_SetString!}.

\begin{center}\rule{0.5\linewidth}{0.5pt}\end{center}

\hypertarget{lazy-iterators-dictionary-view-objects}{%
\subsection{7.4 Lazy Iterators \& Dictionary View
Objects}\label{lazy-iterators-dictionary-view-objects}}

\hypertarget{lazy-conversions-rangeobject-memory-layout}{%
\subsubsection{\texorpdfstring{1. Lazy Conversions \&
\texttt{rangeobject} Memory
Layout}{1. Lazy Conversions \& rangeobject Memory Layout}}\label{lazy-conversions-rangeobject-memory-layout}}

Python 3.0 replaced list-producing functions with lazy iterators. The
\passthrough{\lstinline!range()!} built-in returns a
\passthrough{\lstinline!rangeobject!} defined in
\passthrough{\lstinline!Objects/rangeobject.c!}:

\begin{lstlisting}[language=C]
typedef struct {
    PyObject_HEAD
    PyObject *start;
    PyObject *stop;
    PyObject *step;
    PyObject *length;
} rangeobject;
\end{lstlisting}

Because the sequence elements are not pre-allocated, a
\passthrough{\lstinline!rangeobject!} consumes \(O(1)\) memory. To
compute the value at a specific index \(i\), the type's sequence method
(\passthrough{\lstinline!range\_item!}) performs a direct arithmetic
step calculation: \[\text{value} = \text{start} + i \times \text{step}\]
This math enables \(O(1)\) index access and \(O(1)\) containment checks
for numeric targets:
\[\text{remainder} = (\text{target} - \text{start}) \pmod{\text{step}}\]
If \(\text{remainder} == 0\) and
\(\text{start} \le \text{target} < \text{stop}\) (or reverse for
negative steps), the value is in the range.

\hypertarget{dictionary-view-objects}{%
\subsubsection{2. Dictionary View
Objects}\label{dictionary-view-objects}}

Methods like \passthrough{\lstinline!keys()!},
\passthrough{\lstinline!values()!}, and
\passthrough{\lstinline!items()!} return dynamic views containing a
direct reference back to the parent dictionary, instead of copying
elements into a new list. The underlying struct is
\passthrough{\lstinline!PyDictViewObject!}:

\begin{lstlisting}[language=C]
typedef struct {
    PyObject_HEAD
    PyDictObject *dv_dict;      /* Pointer to parent dictionary struct */
} PyDictViewObject;
\end{lstlisting}
